\documentclass[aspectratio=169,usenames,dvipsnames]{beamer}
\beamertemplatenavigationsymbolsempty
% \documentclass{beamer}

% UNCOMMENT FOR HANDOUTS
% \usepackage{pgfpages}
% \pgfpagesuselayout{2 on 1}[landscape,letterpaper,border shrink=5mm]
% \setbeameroption{show notes}

\mode<presentation>
\usetheme{Boadilla}
\usecolortheme{spruce}
\usepackage{../../LizStyle,ulem}
\usepackage{color}
\usepackage[color,matrix,arrow]{xy}
% \input xy
% \xyoption{all}
\usepackage{tikz}
\usepackage{tikz-cd}
\tikzset{commutative diagrams/diagrams={ampersand replacement=\&}}

\AtBeginSection{\frame{\sectionpage}}


% ==========Command for putting comments and removing them=======
% Visible:
\newcommand{\InstructorNotes}[1]{\textit{\color{orange}#1}}
\newcommand{\InstructorList}[1]{
\begin{itemize}
	\color{orange}
	#1
\end{itemize}
}
% Removed:
% \renewcommand{\InstructorNotes}[1]{}
% \renewcommand{\InstructorList}[1]{}


\newcommand{\bvfill}{\vskip0pt plus 1filll}

\usepackage{multimedia}
\usepackage{graphbox}

% \usepackage{algpseudocode}


\newcommand{\Reeb}{\textbf{Reeb}}
\newcommand{\Rtopc}{\textbf{Rtopc}}
% \newcommand{\RR}{\mathcal{R}}

\newcommand{\Set}{\textbf{Set}}
\newcommand{\Vect}{\textbf{Vect}}
\newcommand{\Top}{\textbf{Top}}
\newcommand{\Int}{\textbf{Int}}


\graphicspath{ {Figures/} {../Figures/}{../../Figures/} }



\begin{document}
%\pdfoptionpdfminorversion 5
%\logo{\includegraphics[height=1.5cm]{dukeshield.pdf}}
\title[]{Ch 5.1.3-4: $k$-Fold Cross-Validation}
\subtitle{Lecture 13 - CMSE 381}
\date{Mon, Sep 29, 2025}



\institute[MSU-CMSE]{Michigan State University\\
			::\\
          Dept of Computational Mathematics, Science \& Engineering}
\author[Dr.~Zhang]{Prof.~Mengsen Zhang}

\begin{frame}
 \titlepage


\end{frame}
%--------------------------
\begin{frame}[c]{Announcements}

\begin{columns}
\column{.45\textwidth}
\centering

	\textbf{Last time:}
	\begin{itemize}
		\item Validation Set 
		\item LOOCV
	\end{itemize}

% \textbf{Announcements:}
% 	\begin{itemize}
% 		\item Exam 1 grades
% 		\item HW \#4 Posted. 
% 		\begin{itemize}
% 			\item Changed Deadline! Due Wednesday Oct 9. 
% 		\end{itemize}
% 	\end{itemize}
\column{.45\textwidth}
\centering
\includegraphics[width = \textwidth]{Schedule-F25-2.png}
\end{columns}

\end{frame}
%--- Next Frame ---%
\begin{frame}[c]{Covered in this lecture }

\begin{columns}
\column{.4\textwidth}
\centering
\begin{itemize}
	\item$k$-fold CV
\end{itemize}
\column{.4\textwidth}
\centering

\end{columns}

\end{frame}
%--- Next Frame ---%
%=================
\section{Last time}
%=================

%--- Next Frame ---%
\begin{frame}[t]{Validation set approach}

\includegraphics[width = \textwidth]{CV-subset0}

\begin{columns}
\column{.45\textwidth}
\centering
\begin{itemize}
	\item Divide randomly into two parts:
	\begin{itemize}
		\item Training set
		\item Validation/Hold-out/Testing set
	\end{itemize}
	\item Fit model on training set
	\item Use fitted model to predict response for observations in the test set
	\item Evaluate quality (e.g.~MSE)
\end{itemize}
\column{.45\textwidth}
\centering
\InstructorList{
\item This example has 30 data points
\item Keep out $30\% = 9 $ for testing in yellow
}
\end{columns}

\end{frame}
%--- Next Frame ---%
\begin{frame}[c]{Problems}
\begin{columns}
\column{.3\textwidth}
\centering
\includegraphics[width = \textwidth]{Fig5-2b}
Ex.~Predict \texttt{mpg} using \texttt{horsepower}
\column{.3\textwidth}

\includegraphics[width = \textwidth]{CV-subset0}
\includegraphics[width = \textwidth]{CV-subset1}
\includegraphics[width = \textwidth]{CV-subset2}
\includegraphics[width = \textwidth]{CV-subset3}
\includegraphics[width = \textwidth]{CV-subset4}

\centering

\column{.3\textwidth}

\begin{itemize}
\item Highly variable results, no consensus about the error
\item Tends to overestimate test error rate
\end{itemize}

\end{columns}

\end{frame}
%--- Next Frame ---%
\begin{frame}[c]{Leave One Out CV (LOOCV)}

\begin{columns}
\column{.45\textwidth}
\centering
\begin{itemize}
	\item Remove $(x_1,y_1)$ for testing.
	\item Train the model on $n-1$ points: $\{(x_2,y_2), \cdots, (x_n,y_n)\}$
	\item Calculate $\mathrm{MSE}_1 = (y_1-\hat y_1)^2$
	\item[]
	\item Remove $(x_2,y_2)$ for testing.
	\item Train the model on $n-1$ points: $\{(x_1,y_1),(x_3,y_3), \cdots, (x_n,y_n)\}$
	\item Calculate $\mathrm{MSE}_2 = (y_2-\hat y_2)^2$
	\item []
	\item Rinse and repeat
\end{itemize}
\column{.45\textwidth}
\centering

\includegraphics[width = \textwidth]{LOOCV-subset0.png}

\includegraphics[width = \textwidth]{LOOCV-subset1.png}

\includegraphics[width = \textwidth]{LOOCV-subset2.png}

\includegraphics[width = \textwidth]{LOOCV-subset3.png}

\includegraphics[width = \textwidth]{LOOCV-subset4.png}
$\vdots$

\includegraphics[width = \textwidth]{LOOCV-subset29.png}
Return the score:
\begin{equation*}
	CV_{(n)} = \frac{1}{n}\sum_{i=1}^n \mathrm{MSE}_i
\end{equation*}
\end{columns}

\end{frame}
%--- Next Frame ---%
\begin{frame}[c]{Pros and Cons}

\begin{columns}
\column{.35\textwidth}
\centering
% \includegraphics[width = \textwidth]{Fig5-2b}
\includegraphics[width = \textwidth]{Fig5-4a}
\column{.45\textwidth}
\centering
\begin{itemize}
	\item No variance
	\item Higher computation cost
\end{itemize}
\end{columns}

\end{frame}
%--- Next Frame ---%
% \begin{frame}[t]{Speeding up LOOCV}

% 	{\color{red}Warning:} This only works for least squares linear or polynomial regression.


% 	\begin{columns}
% 	\column{.45\textwidth}
% 	\centering

% \begin{equation*}
% 	h_i = \frac{1}{n} + \frac{(x_i - \overline x)^2}{\sum_{j=1}^n (x_j - \overline x)^2}
% \end{equation*}

% \InstructorNotes{Leverage}
% 	\column{.45\textwidth}
% 	\centering

% 	\begin{equation*}
% 		\frac{1}{n}\sum_{i=1}^n \mathrm{MSE}_i = CV_{(n)} = \frac{1}{n}\sum_{i=1}^n \left( \frac{y_i - \hat y_i}{1-h_i} \right)^2
% 	\end{equation*}
% 	\InstructorList{
% 	\item This was part of their friday homework
% 	\item Problem is this only works in a VERY limited setting
% 	}
% 	\end{columns}



% \end{frame}
%--- Next Frame ---%
%==========
\section{$k$-Fold CV}
%==========
\begin{frame}[c]{The idea}

\begin{columns}
% \column{.45\textwidth}
% \centering

\column{.75\textwidth}
\centering

\includegraphics[width = \textwidth]{CV-Kfold-subset0}

\includegraphics[width = \textwidth]{CV-Kfold-subset1}

\includegraphics[width = \textwidth]{CV-Kfold-subset2}

\includegraphics[width = \textwidth]{CV-Kfold-subset3}

\includegraphics[width = \textwidth]{CV-Kfold-subset4}
\end{columns}

\InstructorList{
\item $K=5$ in this example
\item Randomly sort, split up into $K$ subsets
\item If $K=n$, then we get LOOCV
}



\end{frame}
%--- Next Frame ---%
\begin{frame}[c]{Mathy version}

\begin{columns}
\column{.45\textwidth}
\centering
\begin{itemize}
	\item Randomly split data into $k$-groups (folds)
	\item Approximately equal sized. For the sake of notation, say each set has $\ell$ points
	\item []
	\item Remove $i$th fold $U_i$ and reserve for testing.
	\item Train the model on remaining points
	\item Calculate $\mathrm{MSE}_i = \frac{1}{\ell} \sum_{(x_j,y_j) \in U_i}(y_j-\hat y_j)^2$
	\item[]
	\item Rinse and repeat
\end{itemize}
\column{.45\textwidth}
\centering
Return
\begin{equation*}
	CV_{(k)} = \frac{1}{k} \sum_{i=1}^k \mathrm{MSE}_i
\end{equation*}
Test your understanding: \href{https://PollEv.com/multiple_choice_polls/upjTEpf6lvhWuTU3a3rep/respond}{PollEv}
\end{columns}

\end{frame}
%--- Next Frame ---%
{
\setbeamercolor{background canvas}{bg=Apricot!25}
\begin{frame}[t]{By hand first!}
There are 10 students in the class, and we have data points for each. They have already been randomly permuted below. 
Write down the training/testing sets for a $3$-fold CV
\begin{columns}
\column{.24\textwidth}
\centering

\begin{itemize}
	\item Damien 
	\item Alice
	\item Greta 
	\item Jasmin
	\item Benji 
	\item Inigo
	\item Frank 
	\item Carina 
	\item Enrique 
	\item Hubert 
\end{itemize}

\column{.24\textwidth}
\centering
\textbf{Fold 1}

\InstructorList{
  \item Test: DAGJ
  \item Train: BIFCEH
}

\vspace{2in}
\column{.24\textwidth}
\centering
\textbf{Fold 2}
\InstructorList{
  \item Test:BIF 
  \item Train: DAGJCEH
}
\vspace{2in}
\column{.24\textwidth}
\centering
\textbf{Fold 3}
\InstructorList{
  \item Test: CEH
  \item Train: DAGJBIF
}
\vspace{2in}
\end{columns}
\end{frame}
}
%--- Next Frame ---%
{
\setbeamercolor{background canvas}{bg=Apricot!25}
\begin{frame}[t]{Coding - Building $k$-fold CV}
\begin{columns}
\column{.45\textwidth}
\centering
\column{.45\textwidth}
\centering
\end{columns}
\end{frame}
}
%--- Next Frame ---%
\begin{frame}[t]{Pros and Cons}

\begin{columns}
\column{.45\textwidth}
\centering
\textbf{Pros:}
\InstructorList{
\item Cheaper computation with $k = 5$ or $10$ instead of $n$
}
\column{.45\textwidth}
\centering
\textbf{Cons:}
\InstructorList{
\item Unlike LOOCV, there is randomness, but it has lower variability  than validation

}

\end{columns}

\end{frame}
%--- Next Frame ---%
\begin{frame}[c]{Comparison}

\begin{columns}
\column{.3\textwidth}
\centering
% \includegraphics[width = \textwidth, align = c]{Fig5-2a}
% \includegraphics[width = \textwidth, align = c]{Fig5-2b}
% \column{.3\textwidth}
% \centering

\includegraphics[width = \textwidth, align = c]{Fig5-4a}
\column{.3\textwidth}
\centering
\includegraphics[width = \textwidth, align = c]{Fig5-4b}
\column{.3\textwidth}
\centering 
\textbf{\tiny Validation set}\\
\includegraphics[width = \textwidth]{Fig5-2b}
\end{columns}

\end{frame}
%--- Next Frame ---%
% \begin{frame}[t]{Comparison with simulated data: Ex 3}

% \begin{columns}
% \column{.35\textwidth}
% \centering
% \includegraphics[width = \textwidth, align = c, trim = 0 0 700 0, clip]{Fig2-11}
% \column{.3\textwidth}
% \centering
% \includegraphics[width = \textwidth, align = c]{Fig5-6c}
% \column{.3\textwidth}
% \InstructorNotes{}
% \InstructorList{
%   \item Blue: true test MSE is shown in   blue
%   \item Black dashed:  the LOOCV estimate 
%   \item Orange: 10-fold CV 
%   \item The crosses indicate the minimum of each of the   MSE curves.
%   \item Nearly the same
% }


% \end{columns}
% \InstructorNotes{}
% \end{frame}
% %--- Next Frame ---%
% \begin{frame}[t]{Comparison with simulated data: Ex 1}

% \begin{columns}
% \column{.35\textwidth}
% \centering
% \includegraphics[width = \textwidth, align = c, trim = 0 0 590 0, clip]{Fig2-9}
% \column{.32\textwidth}
% \centering
% \includegraphics[width = \textwidth, align = c]{Fig5-6a}

% \column{.3\textwidth}
% \InstructorList{
%   \item Blue: true test MSE is shown in   blue
%   \item Black dashed:  the LOOCV estimate 
%   \item Orange: 10-fold CV 
%   \item The crosses indicate the minimum of each of the   MSE curves.
%   \item \textbf{Underestmates the true test MSE}
% }
% \end{columns}
% \end{frame}
% %--- Next Frame ---%
% \begin{frame}[t]{Comparison with simulated data: Ex 2}

% \begin{columns}
% \column{.35\textwidth}
% \centering
% \includegraphics[width = \textwidth, align = c, trim = 0 0 600 0, clip]{Fig2-10}
% \column{.29\textwidth}
% \centering
% \includegraphics[width = \textwidth, align = c]{Fig5-6b}
% \column{.29\textwidth}
% \InstructorNotes{Close for lower flexibility, then overestimates at higher ranges}

% \end{columns}
% \end{frame}
% %--- Next Frame ---%
% \begin{frame}[c]{Takeaways from the examples}

% \begin{columns}
% \column{.45\textwidth}
% \centering
% \InstructorList{
% \item In all plots, LOOCV and 10-fold CV are similar
% \item Sometimes really only care about the minimum point in the test MSE because want to get the right degree of freedom for a model. All examples have approximately same x-coord for that

% }
% \column{.45\textwidth}
% \centering

% \end{columns}

% \end{frame}
% %--- Next Frame ---%
% \begin{frame}[t]{Bias-Variance Tradeoff: Bias}
% 	\begin{equation*}
% 		E(y_0 - \hat f (x_0))^2 = \Var(\hat f (x_0))
% 			+ \left[\mathrm{Bias}(\hat f(x_0)) \right] ^2
% 			+ \Var(\e)
% 	\end{equation*}
% \begin{columns}
% \column{.45\textwidth}
% \centering
% % \textbf{Bias}
% \InstructorList{
% \item 1.... Validation set overestimates test error b/c used small subset of data
% \item 3.... $k$-fold gives medium level of bias b/c training set has approximately $(k-1)n/k$ observations
% \item 2....LOOCV gives approximately unbiased estimate since uses almost all data every time
% }
% \column{.45\textwidth}
% \centering
% \InstructorList{
% \item sooooooooo LOOCV better than $k$-fold just for bias
% }
% \end{columns}

% \end{frame}
% %--- Next Frame ---%
% \begin{frame}[t]{Bias-Variance Tradeoff: Variance}
% 	\begin{equation*}
% 		E(y_0 - \hat f (x_0))^2 = \Var(\hat f (x_0))
% 			+ \left[\mathrm{Bias}(\hat f(x_0)) \right] ^2
% 			+ \Var(\e)
% 	\end{equation*}
% \begin{columns}
% \column{.45\textwidth}
% \centering

% % \textbf{Variance}
% \InstructorList{
% \item LOOCV higher variance:
% \InstructorList{
% 	\item avg the outputs of n fitted models, but all on almost identical observations
% 	\item high correlation with each other
% }
% \item $k$-fold
% \InstructorList{
% \item $k$ fitte models somewhat less correlated with each other
% \item mean of many highly correlated quantities has higher variance than those not corrleated
% }
% }

% \column{.45\textwidth}
% \centering
% \InstructorList{
% \item Empirically, $k=5$ or 10 has been shown to be a happy medium
% }
% \end{columns}

% \end{frame}
% %--- Next Frame ---%
% {
% \setbeamercolor{background canvas}{bg=Apricot!25}
% \begin{frame}[t]{Do the remainder of the coding}
% \begin{columns}
% \column{.45\textwidth}
% \centering
% \column{.45\textwidth}
% \centering
% \end{columns}
% \end{frame}
% }
% %--- Next Frame ---%
% %----------------
\begin{frame}[t]{Next time}


	\includegraphics[align = c, width = .6\textwidth]{Schedule-F25-2.png}
\begin{columns}
\column{.45\textwidth}
\centering


\column{.45\textwidth}

\end{columns}
\end{frame}
%--- Next Frame ---%


\end{document}
